% !TEX root = ../../mixed-sklyanin-cubics.tex
% Current result: 2026-11-22 / v59.
\subsection{中央切断の判別形式と比較を始められる開集合}
\label{mcpn:subsec:pencil}

以下は $k\ge2$、$N=3k$、中心変数二つ $U=(u,v)$ の場合である。
前節で一度固定した境界比較と Jacobian 切断から、
式\eqref{euk:eq:coefficients}の多項式 $Q,L,C$ を固定する。
定理\ref{mcp:thm:normal-form}の係数を
\[
 \mathbf t(U)\in\C^N\otimes\C[u,v]_1,\quad
 \alpha(U)\in\C^N\otimes\C[u,v]_2,\quad c(U)\in\C[u,v]_3
\]
と書く。これらを重み斉次式へ代入して得る三次式を
\begin{equation}\label{mcpn:eq:polarized-candidate}
 \begin{split}
 W_{\mathrm{pol}}(\mathbf x,U)
  ={}&G(\mathbf x)+Q(\mathbf t(U))(\mathbf x)\\
    &+L(\mathbf t(U),\alpha(U))(\mathbf x)
      +C(\mathbf t(U),\alpha(U),c(U)),\\
 X_{\mathrm{pol}}={}&V(W_{\mathrm{pol}})\subset\mathbb P^{N+1}
 \end{split}
\end{equation}
と定める。$Q,L,C$ の active 次数はそれぞれ二、一、零であり、
中心変数に関する次数は一、二、三である。
本小節はこの\emph{偏極した可換候補}の幾何を扱う。
多項式を偏極しただけで二中心の Morita 核が得られたとはしない。

一中心の判別多項式を系\ref{euk:cor:discriminant}の $D$ とし、
$\mu=2^N=8^k$ と置く。

\begin{proposition}[中央切断を判定する次数 $3\cdot8^k$ の形式]
\label{mcpn:prop:ray-discriminant}
\[
 \Delta(u,v)=D(\mathbf t(u,v),\alpha(u,v),c(u,v))
\]
は次数 $3\mu$ の二変数斉次形式である。恒等零となる特殊係数も許す。
$\ell=[a:b]\in\mathbb P^1$ に対し
\[
 A_\ell=A/(bu-av)
\]
と書くと、
\[
 \Delta(a,b)\ne0
 \quad\Longleftrightarrow\quad
 A_\ell\text{ の残余線束圏が smooth}
 \quad\Longleftrightarrow\quad
 V\bigl(W_{\mathrm{pol}}(\mathbf x,as,bs)\bigr)
       \subset\mathbb P^N\text{ が滑らか}
\]
である。従って全二中心係数空間の開集合
\[
 \mathcal U_{\mathrm{cut}}
   =\{\Delta\not\equiv0\}
\]
は、少なくとも一つの滑らかな中央切断を持つ係数を正確に表す。
その補集合は $\Delta$ の全 $3\mu+1$ 個の係数の消滅で定義される。
\end{proposition}

\begin{proof}
一中心側で active 変数を $\mathbf x\mapsto z\mathbf x$ と変えると、
係数の重み $(1,2,3)$ により臨界値も定数項も $z^3$ 倍される。
従って $z\ne0$ について
\[
 D(z\mathbf t,z^2\alpha,z^3c)=z^{3\mu}D(\mathbf t,\alpha,c).
\]
これは全係数の多項式恒等式なので $z=0$ にも成立する。
代入した $\Delta$ の斉次性が従う。
正規形の中心持上げは普遍多項式の底変更であるから、
$(u,v)=(as,bs)$ の代入は、一中心族の係数
$(\mathbf t(a,b),\alpha(a,b),c(a,b))$ と正確に一致する。
系\ref{euk:cor:discriminant}が表示した同値を与える。
非零斉次形式は $\mathbb P^1$ 全体では消えないので、最後の主張も従う。
\end{proof}

\begin{theorem}[一次形式を固定した一般係数の中央 pencil]
\label{mcpn:thm:lefschetz-pencil}
任意の $\mathbf t(U)$ を固定する。
二次形式 $\alpha(U)$ と三次形式 $c(U)$ の係数空間には非空 Zariski 開集合があり、
そこで $\Delta$ は正確に $3\mu$ 個の相異なる零点を持つ。
各零点の可換切断は唯一の通常二重点を持つ。
この開集合では $X_{\mathrm{pol}}$ は滑らかで、
\[
 \operatorname{Bl}_{Z}X_{\mathrm{pol}}\longrightarrow\mathbb P^1,
 \qquad Z=V(G)\subset\{u=v=0\},
\]
は $3\mu$ 個の一節点ファイバーを持つ Lefschetz pencil の解消である。
\end{theorem}

\begin{proof}
係数空間を
$\mathcal P=(\C^N\otimes H^0(\mathbb P^1,\OO(2)))
\times H^0(\mathbb P^1,\OO(3))$ とする。
各 $\ell\in\mathbb P^1$ で、値の評価
$(\alpha(U),c(U))\mapsto(\alpha(\ell),c(\ell))$ は全射である。
さらに $c$ の値と一次 jet は独立に指定できる。

固定した $\mathbf t(\ell)$ に対する $D=0$ は、
系\ref{euk:cor:normalization}により既約被約超曲面である。
その特異部分は $(\alpha,c)$ 空間で余次元二以上である。
値の評価の全射性から、$\mathbb P^1\times\mathcal P$ 上で
この特異部分へ入る incidence も余次元二以上となる。
従ってその $\mathcal P$ への像は真の閉部分集合である。
閉性には $\mathbb P^1$ の proper 性を使う。

この部分を除いた $D=0$ の点では、対応する臨界点が唯一で、
Hessian が非退化である。臨界代数をその長さ一の因子と残りに分けると、
norm の微分は
\[
                   \frac{\partial D}{\partial c}\ne0
\]
となる。ここでは $c\mapsto C$ の係数 $\kappa\ne0$ を使う。
$\Delta(\ell)=0$ を固定したまま $c$ の一次 jet を変えれば、
$d\Delta(\ell)$ を任意に変えられる。
従って $\Delta=d\Delta=0$ の incidence も、この領域では余次元二である。
前の特異部分を含めた閉集合の全次元は高々 $\dim\mathcal P-1$ なので、
その像を除いた非空開集合では全ての零点が単純となる。
次数 $3\mu$ から零点数を得る。

$Z$ 上では $G$ の偏微分が同時に消えないため、
$X_{\mathrm{pol}}$ と全ての切断は滑らかである。
$Z$ の外では中心方向の局所座標を $z$ とし、三次式を
$F(\mathbf x,z)=0$ と書く。切断が滑らかなら total space も滑らかである。
切断の唯一の Morse 点 $p$ では、上の長さ一因子による norm の局所表示から
\[
 d\Delta/dz\ne0\quad\Longrightarrow\quad
                       \partial F/\partial z(p)\ne0.
\]
従ってそこでも total space は滑らかである。
陰関数定理で $z$ を解くと、その Hessian は非退化なので
射影の臨界点は通常の Lefschetz 型となる。

$(u,v)$ は $X_{\mathrm{pol}}$ 上の正則列であり、
その零集合は滑らかな余次元二の $Z$ である。
従って pencil の graph はこの ideal の blow-up に一致する。
方向 $[a:b]$ のファイバーは
$V(W_{\mathrm{pol}}(\mathbf x,as,bs))$ 自身である。
各切断内では $Z$ は Cartier 因子なので余分な blow-up は生じない。
\end{proof}

次数 $3\cdot2^N$ 自体は、滑らかな $N$ 次元三次の双対次数という既知の数である
\cite[Section~7.1, Example~7.4]{CMOTevelevDuality}。
上の証明で確認したのは、$G$ と任意の $\mathbf t(U)$ を固定し、
制限された $Q$ の族の中でも、$\alpha(U),c(U)$ の変化だけで
この一般位置が得られることである。

\begin{example}[六 active 変数における $192$ 個の単純臨界方向]
\label{mcpn:ex:fermat-pencil}
$G=\sum_{j=0}^{5}x_j^3$、$b_j=2^j$ として
\[
 W=\sum_{j=0}^{5}x_j^3-3u^2\sum_{j=0}^{5}b_j^2x_j+v^3
\]
を取る。$u\ne0$ の臨界点は $x_j=\epsilon_jb_ju$、$\epsilon_j\in\{\pm1\}$ であり、
中央切断の判別形式は
\[
 \Delta(u,v)=\prod_{\epsilon\in\{\pm1\}^6}
          \left(v^3-2\sum_{j=0}^{5}\epsilon_j8^j\,u^3\right)
\]
である。符号付き和は互いに異なり、どれも零ではない。
従って $64$ 個の三次因子は互いに素で、各因子の三根は相異なる。
全 $192$ 方向が単純である。Hessian は
$\operatorname{diag}(6\epsilon_jb_ju)$ なので全臨界点が Morse であり、
$u=0$ のファイバーも滑らかである。
この有限の恒等式と jet の評価は
\texttt{computations/verify\_multicentral\_pencil.py}で再現する。
これは一般の比較係数の数値計算や、二中心核の検算ではない。
\end{example}

\begin{remark}[切断の判別と全対象の比較の区別]
\label{mcpn:rem:scope}
$\mathcal U_{\mathrm{cut}}$ 上で選んだ良い切断は、次の逐次比較の出発点になる。
そこで構成する標的と $X_{\mathrm{pol}}$ が同じであることは未証明である。
特に $X_{\mathrm{pol}}$ の smoothness を、そのまま二中心の非可換対象の
smoothness と同一視しない。
$\Delta\equiv0$ の部分に滑らかな非可換対象が存在しないことも、ここでは証明していない。
一般の非可換 Bertini を適用するには追加の確認が要る。
例えば\cite[Proposition~2.1]{CMORSVBertini}は可換有限型環上の有限性を仮定し、
同論文の Introduction は smooth dg 圏の層への拡張を証明していないと明記する。
一般の多中心 PBW 族にその有限中心仮定を無条件で加えることはしない。
\end{remark}
